\documentclass[11pt]{article}
\usepackage[margin=1in]{geometry}
\usepackage{amsmath,amssymb,amsthm}
\usepackage{array}
\usepackage{parskip}
% The conjecture is quoted bilingually; load a CJK font for the Chinese line.
\usepackage{xeCJK}
\setCJKmainfont{Songti SC}

\newtheorem{theorem}{Theorem}
\newtheorem{lemma}[theorem]{Lemma}
\newtheorem{corollary}[theorem]{Corollary}
\newtheorem{proposition}[theorem]{Proposition}
\theoremstyle{definition}
\newtheorem{definition}[theorem]{Definition}
\theoremstyle{remark}
\newtheorem{remark}[theorem]{Remark}

\newcommand{\diam}{\operatorname{diam}}
\DeclareMathOperator{\dist}{dist}

\title{Disproof of conjecture \texttt{00000001678}}
\author{earthking11}
\date{2026-09-13}

\begin{document}
\maketitle

\section*{Verdict: FALSE}

We disprove conjecture \texttt{00000001678} as stated. Its first clause --- that
the difference between the zero forcing number and the path cover number on
trees takes values in $\{0,1\}$ --- is not the point of failure. The failure is
the \emph{``exactly''} clause: the asserted classification of the difference-one
trees. In fact, by the AIM theorem, $Z(T)=P(T)$ for \emph{every} tree $T$, so the
difference is always $0$ and the difference-one class is \emph{empty}, while the
class of odd-diameter trees admitting a perfect matching is non-empty. The
smallest witnesses are the paths $P_2, P_4, P_6, \dots$; the case $T=P_4$ is
worked out below and formalised in Lean~4.

\section{The conjecture}

Quoted verbatim from \texttt{conjectures/00000001678.md}:

\begin{quote}
\textbf{English.} Conjecture: The difference between the zero forcing number and
the path cover number on trees takes values in $\{0, 1\}$, with an explicit tree
classification (zero-forcing versus path-cover tree classification); and the
trees with difference one are exactly the odd-diameter trees admitting a perfect
matching.

\textbf{中文。} 猜想：零力数(zero forcing number)与路径覆盖数之差在树上取值
$\{0, 1\}$ 且树分类显式($Z-$路径覆盖树分类)；且差为一的树恰为含完美匹配的奇直径树。
\end{quote}

\section{Definitions}

Throughout, $G=(V,E)$ is a finite simple graph and, later, a tree.

\begin{definition}[zero forcing number; colour-change rule]
Colour every vertex of $G$ either \emph{black} or \emph{white}. The
\emph{colour-change rule} is: if a black vertex $u$ has \emph{exactly one} white
neighbour $v$, then $v$ is turned black; one says that $u$ \emph{forces} $v$. A
\emph{zero forcing set} is a set $S\subseteq V$ of vertices such that, starting
with exactly the vertices of $S$ black and all others white and repeatedly
applying the colour-change rule, eventually every vertex is black. The
\emph{zero forcing number} $Z(G)$ is the minimum cardinality of a zero forcing
set.
\end{definition}

The iteration terminates because the number of black vertices is non-decreasing
and bounded by $|V|$; the closure of $S$ is well defined.

\begin{definition}[path cover number]
A \emph{path cover} of $G$ is a set of pairwise vertex-disjoint paths whose
vertex sets cover $V$. A path cover is \emph{induced} if each path is an induced
subgraph of $G$. The \emph{path cover number} $P(G)$ is the minimum number of
paths in a path cover of $G$.

For a tree the two conventions coincide: a path subgraph of a tree is
automatically induced, since a tree has no chords (any additional edge between
two vertices of the path would close a cycle).
\end{definition}

\begin{definition}[diameter and perfect matching]
The \emph{diameter} $\diam(G)$ is the maximum graph distance between two
vertices of $G$. A \emph{perfect matching} of $G$ is a set of pairwise disjoint
edges whose endpoints together cover every vertex of $G$.
\end{definition}

The two classes compared by the conjecture are therefore
\[
\begin{aligned}
\mathcal{D}\; &=\;\{\,T \text{ a tree} : Z(T)-P(T)=1\,\}, \\
\mathcal{R}\; &=\;\{\,T \text{ a tree} : \diam(T)\text{ odd and } T \text{ has a perfect matching}\,\},
\end{aligned}
\]
and the ``exactly'' clause asserts $\mathcal{D}=\mathcal{R}$.

\section{The AIM theorem: the difference is always zero}

The following is the relevant theorem from the AIM minimum-rank / zero-forcing
programme.

\begin{theorem}[AIM; Barioli--Fallat--Hogben et al., 2008]
\label{thm:aim}
For every tree $T$, the zero forcing number and the path cover number agree:
\[
Z(T) \;=\; P(T).
\]
\end{theorem}

\begin{proof}[Reference]
F.~Barioli, S.~Fallat, L.~Hogben, et al., \emph{Zero forcing sets and the
minimum rank of graphs}, Linear Algebra and its Applications \textbf{428}
(2008) 1628--1648. The inequality $Z(T)\ge P(T)$ holds for every graph, and
equality for trees is the content of the cited theorem.
\end{proof}

\begin{corollary}\label{cor:empty}
For every tree $T$, $Z(T)-P(T)=0$. Hence the difference-one class $\mathcal{D}$
is empty.
\end{corollary}

\begin{corollary}\label{cor:rhs-nonempty}
The class $\mathcal{R}$ is non-empty: for instance every path $P_{2k}$ has odd
diameter $2k-1$ and the perfect matching $\{v_1v_2,\,v_3v_4,\,\dots,\,v_{2k-1}v_{2k}\}$.
\end{corollary}

Since $\mathcal{D}=\varnothing$ while $\mathcal{R}\neq\varnothing$ by the two
corollaries, $\mathcal{D}=\mathcal{R}$ is impossible. The remainder of the paper
makes this concrete by evaluating a single smallest witness.

\section{The witness $T=P_4$}

Let $P_4$ have vertices $v_1,v_2,v_3,v_4$ and edges $v_1v_2$, $v_2v_3$,
$v_3v_4$.

\begin{lemma}[$Z(P_4)=1$]\label{lem:z}
\end{lemma}
\begin{proof}
\emph{Upper bound.} Take the zero forcing set $S=\{v_1\}$. Vertex $v_1$ is black
and has the unique white neighbour $v_2$, so $v_1$ forces $v_2$. Now $v_2$ is
black; its neighbours are $v_1$ (black) and $v_3$ (white), so $v_2$ forces
$v_3$. Now $v_3$ is black; its neighbours are $v_2$ (black) and $v_4$ (white),
so $v_3$ forces $v_4$. All four vertices are black. Thus $Z(P_4)\le 1$.

\emph{Lower bound.} The empty set has no black vertex, so the colour-change rule
never applies and no vertex is ever forced; since $|V|=4>0$ the empty set is not
a zero forcing set. Hence $Z(P_4)\ge 1$. Therefore $Z(P_4)=1$.
\end{proof}

The explicit forcing sequence is
\[
\{v_1\}\;\xrightarrow{v_1\to v_2}\;\{v_1,v_2\}\;
\xrightarrow{v_2\to v_3}\;\{v_1,v_2,v_3\}\;
\xrightarrow{v_3\to v_4}\;\{v_1,v_2,v_3,v_4\}.
\]

\begin{lemma}[$P(P_4)=1$]\label{lem:p}
\end{lemma}
\begin{proof}
\emph{Upper bound.} The single path $v_1v_2v_3v_4$ covers all four vertices and
is a subgraph of $P_4$; as noted above, in a tree every path subgraph is
induced, so it is an admissible (induced) path cover with one path. Hence
$P(P_4)\le 1$.

\emph{Lower bound.} Every path cover has at least one path because $V\neq
\varnothing$. Hence $P(P_4)\ge 1$. Therefore $P(P_4)=1$.
\end{proof}

\begin{theorem}[the difference on $P_4$ is zero]\label{thm:diff}
\[
Z(P_4)-P(P_4)\;=\;1-1\;=\;0\;\neq\;1 .
\]
In particular $P_4\notin\mathcal{D}$.
\end{theorem}
\begin{proof}
Immediate from Lemma~\ref{lem:z} and Lemma~\ref{lem:p}.
\end{proof}

\begin{lemma}[$P_4$ is an odd-diameter tree with a perfect matching]\label{lem:rhs}
$\diam(P_4)=3$, which is odd, and $\{v_1v_2,\,v_3v_4\}$ is a perfect matching
of $P_4$. Consequently $P_4\in\mathcal{R}$.
\end{lemma}
\begin{proof}
Along a path the graph distance is the index difference, so
$\dist(v_i,v_j)=|i-j|$. The largest value is $\dist(v_1,v_4)=3$; thus
$\diam(P_4)=3$, an odd number. The edges $v_1v_2$ and $v_3v_4$ are disjoint and
their endpoints $\{v_1,v_2\}$ and $\{v_3,v_4\}$ cover all four vertices, so they
form a perfect matching.
\end{proof}

\begin{theorem}[the ``exactly'' clause is false]\label{thm:main}
Conjecture \texttt{00000001678} is false as stated. Indeed
\[
P_4\in\mathcal{R}\quad\text{but}\quad P_4\notin\mathcal{D},
\]
so $\mathcal{D}\neq\mathcal{R}$ and the assertion that the difference-one trees
are \emph{exactly} the odd-diameter trees admitting a perfect matching fails.
\end{theorem}
\begin{proof}
Lemma~\ref{lem:rhs} gives $P_4\in\mathcal{R}$; Theorem~\ref{thm:diff} gives
$P_4\notin\mathcal{D}$.
\end{proof}

\begin{remark}[even smaller and infinite witnesses]
$P_2$ is an even smaller counterexample: $Z(P_2)=P(P_2)=1$, its diameter is $1$
(odd), and its unique edge is a perfect matching, yet the difference is $0$.
More generally, for every $k\ge 1$ the path $P_{2k}$ has diameter $2k-1$ (odd),
the perfect matching $\{v_1v_2,\dots,v_{2k-1}v_{2k}\}$, and $Z=P=1$, so
\[
P_{2k}\in\mathcal{R}\setminus\mathcal{D}
\qquad\text{for every } k\ge 1 .
\]
The witnesses are therefore infinite in number.
\end{remark}

\begin{remark}[what remains true]
The first clause of the conjecture, that $Z-P\in\{0,1\}$ on trees, is compatible
with the AIM theorem ($Z-P=0$ always), so that clause alone is not refuted. What
is refuted is the classification of the difference-one trees, which is vacuous.
\end{remark}

\section{Computational corroboration}

The accompanying program \texttt{reproduce.py} (Python~3, standard library only)
brute-forces $Z$ over all subsets with the colour-change closure and $P$ by
exhaustive path-cover search. It verifies the table for $P_n$, $n=1,\dots,10$:

\begin{center}
\begin{tabular}{c|r|r|r|r|c|c}
$n$ & $Z(P_n)$ & $P(P_n)$ & $Z-P$ & $\diam$ & odd? & perfect matching? \\
\hline
$1$  & $1$ & $1$ & $0$ & $0$ & no  & no  \\
$2$  & $1$ & $1$ & $0$ & $1$ & yes & yes \\
$3$  & $1$ & $1$ & $0$ & $2$ & no  & no  \\
$4$  & $1$ & $1$ & $0$ & $3$ & yes & yes \\
$5$  & $1$ & $1$ & $0$ & $4$ & no  & no  \\
$6$  & $1$ & $1$ & $0$ & $5$ & yes & yes \\
$7$  & $1$ & $1$ & $0$ & $6$ & no  & no  \\
$8$  & $1$ & $1$ & $0$ & $7$ & yes & yes \\
$9$  & $1$ & $1$ & $0$ & $8$ & no  & no  \\
$10$ & $1$ & $1$ & $0$ & $9$ & yes & yes \\
\end{tabular}
\end{center}

\noindent
In particular every odd-diameter $P_n$ admitting a perfect matching
($n=2,4,6,8,10$) has $Z-P=0$. Enumerating all trees on $1,\dots,10$ vertices
(the $201$ isomorphism classes; Prüfer sequences for $n\le 8$, canonical
rooted trees for $n=9,10$, cross-validated) gives:

\begin{center}
\begin{tabular}{r|r}
trees with $Z-P=1$ & $0$ \\
trees with odd diameter and a perfect matching & $14$ \\
trees with $Z=P$ (AIM) & $201$ \\
\end{tabular}
\end{center}

\noindent
The difference-one class is empty, while the class of odd-diameter trees
admitting a perfect matching is not. The script prints \texttt{PASS} and exits
with status $0$ if and only if these verifications succeed.

\section{Formalisation}

The witness $P_4$ is formalised in Lean~4 in the accompanying \texttt{lean4/}
project (core Lean only, \texttt{import Std}, no Mathlib, no \texttt{sorry}). In
it, vertices are of type \texttt{Fin 4}, adjacency is $i\text{---}(i+1)$, the
colour-change closure is a computable function of the colouring, the zero
forcing number is a minimum over all $16$ subsets, and the path cover number is
a minimum over all $256$ idempotent vertex maps encoding partitions. The main
theorems are

\begin{quote}\ttfamily
theorem singleton\_zero\_forces : isZFS (single 0) = true\\
theorem zeroForcingNumber\_P4 : zeroForcingNumber = 1\\
theorem pathCoverNumber\_P4 : pathCoverNumber = 1\\
theorem diameter\_P4 : diameter = 3\\
theorem diameter\_P4\_odd : isOdd diameter = true\\
theorem P4\_has\_perfect\_matching : hasPerfectMatching = true\\
theorem difference\_P4 : zMinusP = 0\\
theorem difference\_P4\_ne\_one : zMinusP \(\ne\) 1\\
theorem conjecture\_00000001678\_refuted : \dots
\end{quote}

\noindent
Every theorem is discharged by computation (\texttt{decide}); \texttt{Check.lean}
prints \texttt{\#print axioms} for them and shows that no proof uses
\texttt{sorryAx} (only \texttt{propext}, from the standard library, appears).

\section{Files}

\begin{center}
\begin{tabular}{l|p{11cm}}
File & Description \\
\hline
\texttt{README.md} & Verdict, conjecture quote, definitions, the $P_4$ computation, the $P_{2k}$ family, the AIM context, file list, reproduction commands, submission status. \\
\texttt{main.tex} & This LaTeX source (standalone \texttt{article}, builds with \texttt{tectonic main.tex}). \\
\texttt{build/main.pdf} & The compiled document. \\
\texttt{reproduce.py} & Python~3 brute force: $Z$ by colour-change closure, $P$ by exhaustive path covers, the $P_n$ table, and all trees on $1,\dots,10$ vertices. \\
\texttt{lean4/lean-toolchain} & Pins \texttt{leanprover/lean4:v4.33.1}. \\
\texttt{lean4/lakefile.toml} & Lake configuration; library \texttt{Main}, no dependencies. \\
\texttt{lean4/Main.lean} & Lean~4 formalisation of the witness (core Lean, no \texttt{sorry}). \\
\texttt{lean4/Check.lean} & \texttt{\#print axioms} audit of every theorem. \\
\texttt{lean4/README.md} & Statement table, proof strategy, scope note. \\
\end{tabular}
\end{center}

\section{Reproducing}

\begin{quote}\small\ttfamily
\$ python3 reproduce.py\\
\dots\\
PASS: P\_4 has Z = 1, P = 1, Z - P = 0, odd diameter 3 and a\\
perfect matching\\
\$ tectonic main.tex\\
\$ cd lean4 \&\& lake build\\
\$ cd lean4 \&\& lake env lean Check.lean
\end{quote}

\noindent
The numerical checks are dependency-free. The Lean project uses the toolchain
pinned in its \texttt{lean-toolchain} file and imports only \texttt{Std}.

\section{Status against the submission rules}

The submission rules require the LaTeX source code, a PDF document, and a Lean~4
project.

\begin{itemize}
\item \textbf{LaTeX source} --- present (\texttt{main.tex}), a standalone
  \texttt{article} using only \texttt{geometry}, \texttt{amsmath},
  \texttt{amssymb}, \texttt{amsthm}, \texttt{array}, and \texttt{parskip};
  it compiles with \texttt{tectonic main.tex} without external figures or
  non-standard packages.
\item \textbf{PDF document} --- \texttt{build/main.pdf}, produced by
  \texttt{tectonic}; it is a rendered copy of \texttt{main.tex} and exists
  non-empty in this submission.
\item \textbf{Lean 4 project} --- present under \texttt{lean4/} with
  \texttt{lean-toolchain}, \texttt{lakefile.toml}, \texttt{Main.lean}, and
  \texttt{Check.lean}. \texttt{lake build} exits $0$, and \texttt{lake env lean
  Check.lean} prints the axioms of each theorem with no \texttt{sorryAx} and no
  Mathlib dependency. The formalisation establishes, by finite computation, that
  $Z(P_4)=1$, $P(P_4)=1$, $\diam(P_4)=3$ is odd, $P_4$ has a perfect matching,
  and $Z(P_4)-P(P_4)=0\neq 1$.
\end{itemize}

\end{document}
